% macros.tex -- notation, environments and helpers for the paper and its
% supplement. Loaded by preamble.tex after all packages (cleveref must
% already be loaded, so that theorem-like environments with a shared counter
% keep their own reference names). See development/labels.md for usage notes.

% ---------------------------------------------------------------------------
% Theorem-like environments: one counter, numbered within sections
% (Theorem 4.6, Lemma 4.1, ...; in the appendices A.1, B.2, ...; in the
% supplement S1.3, S5.1, ...).
% ---------------------------------------------------------------------------
\theoremstyle{plain}
\newtheorem{theorem}{Theorem}[section]
\newtheorem{lemma}[theorem]{Lemma}
\newtheorem{proposition}[theorem]{Proposition}
\newtheorem{corollary}[theorem]{Corollary}
\theoremstyle{definition}
\newtheorem{definition}[theorem]{Definition}
\newtheorem{assumption}[theorem]{Assumption}
\theoremstyle{remark}
\newtheorem{remark}[theorem]{Remark}
\newtheorem{observation}[theorem]{Observation}

% Hypotheses with a fixed tag instead of a number, e.g. (H) and (H0):
%   \begin{hypothesis}{H}[display]  ...  \end{hypothesis}
% prints "Hypothesis H (display)."; \cref gives "Hypothesis H".
\theoremstyle{definition}
\newtheorem{hypothesisinner}{Hypothesis}
\NewDocumentEnvironment{hypothesis}{m o}
  {\renewcommand{\thehypothesisinner}{#1}%
   \IfValueTF{#2}{\hypothesisinner[#2]}{\hypothesisinner}}
  {\endhypothesisinner}
\theoremstyle{plain}

\crefname{theorem}{Theorem}{Theorems}
\crefname{lemma}{Lemma}{Lemmas}
\crefname{proposition}{Proposition}{Propositions}
\crefname{corollary}{Corollary}{Corollaries}
\crefname{definition}{Definition}{Definitions}
\crefname{assumption}{Assumption}{Assumptions}
\crefname{remark}{Remark}{Remarks}
\crefname{observation}{Observation}{Observations}
\crefname{hypothesisinner}{Hypothesis}{Hypotheses}

% ---------------------------------------------------------------------------
% Boxes
% ---------------------------------------------------------------------------
% Certificate box (one per theorem; fields listed in development/labels.md):
%   \begin{certbox}[Certificate for \cref{thm:lnts-opt}]
%     \certitem{Inputs}{...}
%     ...
%   \end{certbox}
\newtcolorbox{certbox}[1][Certificate]{%
  breakable, enhanced jigsaw, colback=black!3, colframe=black!55,
  boxrule=0.5pt, arc=1pt, left=6pt, right=6pt, top=4pt, bottom=4pt,
  fonttitle=\small\bfseries, coltitle=black, colbacktitle=black!12,
  title={#1}, fontupper=\small}
\newcommand{\certitem}[2]{\par\noindent\hangindent=1.5em\textsf{#1:}\enspace #2\par}

% Numbered text box ("Box 1"), e.g. results at a glance, the camshape800
% example, the SCIP defect:
%   \begin{infobox}[label={box:glance}]{Results at a glance} ... \end{infobox}
\newtcolorbox[auto counter, crefname={Box}{Boxes}]{infobox}[2][]{%
  breakable, enhanced jigsaw, colback=blue!3, colframe=blue!40!black,
  boxrule=0.5pt, arc=1pt, left=6pt, right=6pt, top=4pt, bottom=4pt,
  fonttitle=\small\bfseries, coltitle=black, colbacktitle=blue!10,
  title={Box~\thetcbcounter: #2}, fontupper=\small, #1}

% ---------------------------------------------------------------------------
% Instance names: \inst{waterno2_06} sets the MINLPLib name in typewriter
% type; underscores need no escaping. Line breaks are allowed only after
% an underscore or a hyphen (xurl's break-anywhere rule is switched off).
% The name is always upright: in italic text (theorem statements) the slanted
% typewriter font left a visible gap after each underscore.
% ---------------------------------------------------------------------------
\DeclareUrlCommand\instname{\def\UrlFont{\ttfamily\upshape}%
  \def\UrlBreaks{\do\_\do\-}\def\UrlBigBreaks{}\def\UrlNoBreaks{}%
  \Urlmuskip=0mu\relax}
\DeclareRobustCommand{\inst}{\instname}
\pdfstringdefDisableCommands{\def\inst#1{\detokenize{#1}}}

% ---------------------------------------------------------------------------
% Notation (development/outline.md, section 4)
% ---------------------------------------------------------------------------
% Model, feasible set, optimal value
\newcommand{\model}{\mathcal{M}}                  % stored OSIL model M (M_I: \model_I)
\newcommand{\feas}[1][\model]{\mathcal{F}(#1)}    % exactly feasible set F(M)
\newcommand{\vstar}{v^{*}}                        % optimal value: \vstar(\model)
% Certified bound, primal value, gaps, sense
\newcommand{\Lcert}{L}                            % certified (valid) dual bound
\newcommand{\Uprim}{U}                            % upper end of the objective enclosure at an exactly feasible point
\newcommand{\gapabs}{\Delta}                      % absolute gap s(U - L)
\newcommand{\gaprel}{\delta}                      % relative gap Delta / min(|L|,|U|)
\newcommand{\sense}{s}                            % +1 minimize, -1 maximize
% Data readings (a) GAMS text, (b) OSIL text [all claims], (c) binary64 data
\newcommand{\reading}[1]{reading~(#1)}
% Arithmetic tags [E] exact, [I] mpmath interval, [F] own outward-rounded binary64, [L] library hypothesis
\newcommand{\arith}[1]{\textsf{[#1]}}
% Evidence levels \evid{hand}, \evid{stored}, \evid{rerun} (development/terminology.md) and
% trust-base terms T-int, T-fp, T-mp, T-read, T-code, T-thm
\newcommand{\evid}[1]{\textsf{#1}}
\newcommand{\trust}[1]{\textsf{T-#1}}
% Floating point and listing strings
\newcommand{\fl}{\operatorname{fl}}               % round-to-nearest binary64
\newcommand{\uround}{\mathbf{u}}                  % unit roundoff 2^{-53}
\newcommand{\dunit}[1]{u(#1)}                     % unit of the last displayed digit of string #1
\newcommand{\dval}[1]{d(#1)}                      % rational value of the decimal string #1
% KAN objects: network relaxation R and OSIL model without partition-of-unity rows R_P
\newcommand{\Rnet}{\mathcal{R}}
\newcommand{\RP}{\mathcal{R}_{\mathrm{P}}}
% Number sets and operators
\newcommand{\R}{\mathbb{R}}
\newcommand{\Q}{\mathbb{Q}}
\newcommand{\Z}{\mathbb{Z}}
\newcommand{\N}{\mathbb{N}}
\DeclareMathOperator{\proj}{proj}
\DeclareMathOperator{\dist}{dist}
\DeclareMathOperator*{\argmin}{arg\,min}
\DeclareMathOperator*{\argmax}{arg\,max}
% Scientific notation in text or math: \sci{3.1}{-9} -> 3.1 x 10^-9
\newcommand{\sci}[2]{\ensuremath{#1\cdot 10^{#2}}}

% ---------------------------------------------------------------------------
% Placeholders (each use also writes a warning to the log)
% ---------------------------------------------------------------------------
\definecolor{todocolor}{rgb}{0.8,0,0}
\DeclareRobustCommand{\TODO}[1]{%
  \textcolor{todocolor}{\textbf{[TODO: #1]}}%
  \GenericWarning{}{LaTeX Warning: TODO: #1}}
\DeclareRobustCommand{\archiveDOI}{%
  \textcolor{todocolor}{\texttt{[archive DOI]}}%
  \GenericWarning{}{LaTeX Warning: archive DOI placeholder}}
